\documentclass[Total.tex]{subfiles}

\begin{document}
		\chapter{The Space $L^p$ and Spaces of Measures}
	\section{Space $L^p$}
	\begin{Def}
		Let $(X,\mathcal{A},\mu)$ be a measure space, with nonegative measure $\mu$, and $p\in [1,\infty)$. Then denote
		\begin{gather*}
			\mathcal{L}^p(\mu):=\{f:X\rightarrow [-\infty,\infty]||f|^p\mbox{ is }\mu\mbox{-integrable}\}
		\end{gather*}
		
		Define $\mu$-equivalent: The sets $\mathcal{L}^p(\mu)$ are equipped with the equivalent relation, $\mu$-almost everywhere equivalent
		\begin{gather*}
			f\sim g\Leftrightarrow \mu(\{x\in X|f(x)\not=g(x)\})=0
		\end{gather*}
		Then, denote $L^p(\mu)$ as the factor-space of $\mathcal{L}^p(\mu)$ with respect to the equivalence relation. 
		
		The norm structure is specified by
		\begin{gather*}
			\Vert-\Vert_{L^p(X,\mu)}:f\mapsto (\int_X |f|^p d\mu)^{\frac{1}{p}}
		\end{gather*} 
	\end{Def}
	
	Furthermore, with the limit
	\begin{gather*}
		\lim_{p\rightarrow\infty} (\int_X |f|^p d\mu)^\frac{1}{p}
	\end{gather*}
	\subsubsection{Norm $\Vert\cdot\Vert_p$}
	\begin{lemma}
		For all $\lambda\in \RR,f,g\in \mathcal{L}^p(\mu)$, we have
		\begin{gather*}
			\Vert\lambda f\Vert_p=|\lambda|\Vert f\Vert_p\qquad \Vert f+g\Vert_p\leq \Vert f\Vert_p+\Vert g\Vert_g
		\end{gather*}
	\end{lemma}
	\begin{proof}
		Easy to check, by property of integration. 
	\end{proof}
	
	\begin{corollary}
		The function $\Vert\cdot\Vert_p$ is a norm on the space $L^p(\mu)$. 
	\end{corollary}
	\begin{proof}
		Omitted. 
	\end{proof}
	
	\begin{theorem}
		The space $L^p(\mu)$ is compelete, $p\in [1,+\infty]$. 
	\end{theorem}
	\begin{proof}
		Consider the Cauchy sequence 
	\end{proof}
	\subsection{Convolutions}
		\begin{Def}
			Suppsoe $f,g\in L^1(G)$. Define the \textbf{convolution}
			\begin{gather*}
				f*g(x):=\int_G f(y)g(y^{-1}x)d\lambda(y)
			\end{gather*}
			
			Especially, when $G=(\RR^n,+)$, then
			\begin{gather*}
				f*g(x)=\int_{\RR^n} f(y)g(x-y)dy
			\end{gather*}
		\end{Def}
		
		Especially, change of variable $z=x^{-1}y$,
		\begin{gather*}
			=\int_G f(xz)g(z^{-1})d\lambda(z)
		\end{gather*}
		
		\begin{proposition}
			For all $f,g,h\in L^1(G)$
			\begin{gather*}
				内容...
			\end{gather*}
		\end{proposition}
		\begin{proof}
			内容...
		\end{proof}
		
		
		\subsubsection{Convolution Inequalities}
			\begin{theorem}
				(Minkowski with Convolution)
				\begin{gather*}
					\Vert g* f\Vert_{L^p(G)}\leq \Vert g\Vert_{L^1(G)}\Vert f\Vert_{L^p(G)}
				\end{gather*}
			\end{theorem}
			\begin{proof}
				内容...
			\end{proof}
			
			\begin{theorem}
				(Young)
			\end{theorem}
			\begin{proof}
				内容...
			\end{proof}
			
			\begin{theorem}
				(Young)
			\end{theorem}
			\begin{proof}
				内容...
			\end{proof}
		\subsubsection{Approximate Identities}
			\begin{Def}
				内容...
			\end{Def}
			\begin{Def}
				Let $T$ be an operator defined on a linear space of complex-valued measure functions on a measure space
				
				Then, $T$ is called \textbf{linear} if for all $f,g\in T$, and $\lambda\in \CC$, we have
				\begin{gather*}
					内容...
				\end{gather*}
			\end{Def}
	\subsection{Properties}
	
	
	\begin{Def}
		For $0<p<\infty$, the space $L_{loc}^p(\RR^n,|\cdot|)$ or simply $L_{loc}^p(\RR^n)$ is the set of all Lesbegue0mesurable functions $f$ on $\RR^n$, that satisfies
		\begin{gather*}
			内容...
		\end{gather*}
		
		is the \textbf{locally $p$-integrable function} on $\RR^n$. 
	\end{Def}
	
	\begin{theorem}
		(H\"older) Let $0<p,p_1,\dots,p_k\leq \infty$, where $k\geq 2$, and let $f_j$ be 
	\end{theorem}
	
	\begin{theorem}
		(Jensen)
	\end{theorem}
	
	\begin{theorem}
		\begin{itemize}
			\item (Discrete Minkowski)
				\begin{itemize}
					\item 
					\item 
				\end{itemize}
			\item (Integral Inequality)
			
		\end{itemize}
	\end{theorem}
	
	\begin{proof}
		内容...
	\end{proof}
	
	\begin{theorem}
		(Young)
	\end{theorem}
	
	\begin{theorem}
		(Hardy's Inequality)
	\end{theorem}
	
	\subsection{Approximation of $L^p$}
		
		\begin{lemma}
			The set of all simple functions is everywhere dense in every space $L^p(\mu),1\leq p<\infty$. 
		\end{lemma}
		\begin{proof}
			内容...
		\end{proof}
		
		\begin{corollary}
			Let a nonnegative Borel measure $\mu$ on $\RR^n$ be bounded on bounded sets.
			
			Then, the class $C_0^\infty(\RR^n)$ of \underline{smooth functions with bounded support} is everywhere dense in $L^p(\mu),1\leq p<\infty$. 
			
			In particular, the spaces $L^p(\mu),1\leq p<\infty$, are separable.
		\end{corollary}
		\begin{proof}
			内容...
		\end{proof}
		
		\begin{Def}
			Let $\rho$ be a function integrable over $\RR^n$ such that
			\begin{gather*}
				\int_{\RR^n} \rho(x)dx=1
			\end{gather*}
			
			Now, set $\rho_\epsilon(x)=\frac{\rho(\frac{x}{\epsilon})}{\epsilon^n},\epsilon>0$. 
		\end{Def}
		
		\begin{lemma}
			Let $f\in \mathcal{L}^p(\RR^n),1\leq p<\infty$, then the mapping
			\begin{gather*}
				T_f:\RR^n\rightarrow L^p(\RR^n)\qquad T_f(v)(x)=f(x+v)
			\end{gather*}
			
			is continuous and bounded.
		\end{lemma}
		\begin{proof}
			内容...
		\end{proof}
		
		\begin{theorem}
			Let $f\in \mathcal{L}^p(\RR^n)$, $1\leq p<\infty$, then one has
			\begin{gather*}
				\lim_{\epsilon\rightarrow 0}\Vert f*\rho-f\Vert_p=0
			\end{gather*}
		\end{theorem}
		\begin{proof}
			内容...
		\end{proof}
		
		\begin{corollary}
			内容...
		\end{corollary}
		\begin{proof}
			内容...
		\end{proof}
	\subsection{The Weak Topology in $L^p$}
	
	\section{Uniform Integrability}
	\section{Convergence of Measures}
	
	\section{The Hilbert Space $L^2$}
		\begin{Def}
			Let $\mu$ be a measure with values in $[0,+\infty]$.
			
			The space $L^2(\mu)$ is distinguished among other $L^p(\mu)$ by the property that it is 
			\begin{itemize}
				\item Euclidean: With the inner product
				\begin{gather*}
					(f,g)=\int_X fgd\mu
				\end{gather*}
				
				\item or, Hermitian: With the inner product
				\begin{gather*}
					(f,g)=\int_X f\bar{g}d\mu
				\end{gather*}
			\end{itemize}
		\end{Def}
		
		
		Notice: From $|fg|\leq f^2+g^2$, 
		\begin{gather*}
			f,g\in \mathcal{L}^2(\mu)\Rightarrow fg\in \mathcal{L}^1(\mu)
		\end{gather*}
		
		\begin{proposition}
			Let $H_0$ be a closed linear subspace in a Hilbert space.
			
			Then, $H_0^\perp=\{x\in H|x\perp h,\forall h\in H_0\}$ is a closed linear subspace in $H$, and 
			\begin{gather*}
				H=H_0\oplus H_0^\perp
			\end{gather*}
			
			Hence, for every $h\in H$, there exists a unique vector $h_0$, with $h-h_0\in H_0^\perp$. In addition
			\begin{gather*}
				\Vert h-h_0\Vert=\inf\{\Vert h-x\Vert|x\in H_0\}
			\end{gather*}
		\end{proposition}
		\begin{proof}
			内容...
		\end{proof}
		
		\begin{corollary}
			Let $f$ be a continuous linear function on a Hilbert space $H$.
			
			Then, there exists a unique vector $v$ such that
			\begin{gather*}
				内容...
			\end{gather*}
		\end{corollary}
		\begin{proof}
			内容...
		\end{proof}
		
		\begin{corollary}
			\begin{itemize}
				\item 
				There exists a family of mutually orthogonal unit vector $e_\alpha\in L^2(\mu)$ such that: Every element $f$ in $L^2(\mu)$ is the sum of the following series covergent in $L^2(\mu)$:
				\begin{gather*}
					f=\sum_\alpha c_\alpha e_\alpha
				\end{gather*}
				
				where at most countably many coefficients $c_\alpha$ may be nonzero. 
				
				In addition, one has
				\begin{gather*}
					c_\alpha=(f,e_\alpha)\qquad \Vert f\Vert^2=\sum_\alpha |c_\alpha|^2
				\end{gather*}
				
				The family is calleld an \textbf{orthogonal basis} of space $L^2(\mu)$.
				\item 
			\end{itemize}
		\end{corollary}
		\begin{proof}
			内容...
		\end{proof}
		
		\begin{theorem}
			For any orthonormal system $\{\varphi_n\}$ in $L^2(\mu)$, and any sequence $\{c_n\}\in l^2$, the series
			\begin{gather*}
				\sum_{n=1}^\infty c_n\varphi_n\in L^2(\mu)
			\end{gather*}
			
		\end{theorem}
		\begin{proof}
			内容...
		\end{proof}
		
		\begin{theorem}
			Let $\{\varphi_n\}$ be an orthonormal sequence $L^2(\mu)$.
			
			Then, for all $f\in L^2(\mu)$, the following Bessel inequality holds
			\begin{gather*}
				\sum_{n=1}^\infty |(f,\varphi_n)|^2\leq \Vert f\Vert^2_{L^2(\mu)}
			\end{gather*}
			
			In particular, this equality is true if $\{\varphi_n\}$ is an orthonormal basis. 
		\end{theorem}
		\begin{proof}
			内容...
		\end{proof}
		\section{Duality of the Space $L^p$}
	\chapter{Weak $L^p$ Spaces}
	
		
		\section{Definitions and Properties}
		\begin{Def}
			Take the following of distribution function
			\begin{gather*}
				d_f(\alpha):=\mu(\{x\in X||f(x)>\alpha\})\qquad \alpha>0
			\end{gather*}
			
			Let $f:X\rightarrow \CC$ be measurable, with $p\in [1,\infty)$, then define
			\begin{gather*}
				\Vert f\Vert_{L^{p,\infty}}:=\inf\{C>0|d_f(\alpha)\leq (\frac{C}{\lambda})^p,\forall \alpha>0\}=\sup_{\lambda>0} \lambda d_f(\lambda)^{\frac{1}{p}}
			\end{gather*}
		\end{Def}
		
		\subsection{$L^p$ and $L^{p,\infty}$}
		
		\begin{proposition}
			There is embedding 
			\begin{gather*}
				\mathcal{L}^p(X,\mu)\hookrightarrow \mathcal{L}^{p,\infty}(X,\mu)
			\end{gather*}
		\end{proposition}
		\begin{proof}
			内容...
		\end{proof}
		
		\subsection{Covergence in Measure}
	\chapter{Lorentz Spaces}
		\section{Definitions}
			\begin{Def}
				内容...
			\end{Def}
			
			
			\subsection{Duals of Lorentz Spaces}
	\chapter{Interpolation Theory of Lesbegue Space}
		In general, we will develope the interpolation theory in Banach space. Here we just mainly concern 
		\begin{itemize}
			\item 
			\item 
		\end{itemize}
		
		
		\begin{proposition}
			(Interpolation in Lesbegue Space)Let $0<p<q\leq \infty$, and let $f\in L^{p,\infty}(X,\mu)\cap L^{q,\infty}(X,\mu)$, where $X$ is a $\sigma$-finite measure space. 
			
			Then, $f\in L^r(X,\mu)$, for all $p<r<q$ and
			\begin{gather*}
				\Vert f\Vert_{L^r}\leq (\frac{r}{r-p}+\frac{r}{q-r})^{\frac{1}{r}}\Vert f\Vert_{L^{p,\infty}}^{\frac{\frac{1}{r}-\frac{1}{q}}{\frac{1}{p}-\frac{1}{q}}}\Vert f\Vert_{L^{q,\infty}}^{\frac{\frac{1}{p}-\frac{1}{r}}{\frac{1}{p}-\frac{1}{q}}}
			\end{gather*}
		\end{proposition}
		\begin{proof}
			内容...
		\end{proof}
		
		\section{Real Interpolation}
		\section{Complex Interpolation}
		\section{Interpolation of Operators}
			\subsection{Interpolation of Analytic Family of Operators}
\end{document}